\documentclass[11pt]{article}

\usepackage[margin=1.15in]{geometry}
\usepackage{amsmath,amssymb}
\usepackage{booktabs}
\usepackage{array}
\usepackage{longtable}
\usepackage{microtype}
\usepackage{enumitem}
\emergencystretch=3em
\usepackage{hyperref}
\hypersetup{colorlinks=true,linkcolor=black,urlcolor=black,citecolor=black}

\newcommand{\file}[1]{\texttt{#1}}
\newcommand{\lf}[1]{(landscape\_findings~\S#1)}
\newcommand{\sm}[1]{(search.md~\S#1)}
\newcommand{\sfe}[1]{(solver\_features~\S#1)}
\newcommand{\ap}[1]{(approaches~\S#1)}
\newcommand{\sig}{\sigma}
\newcommand{\sione}{\sig_1}
\newcommand{\sitwo}{\sig_2}
\newcommand{\lam}{\lambda}

\title{The eigenvalue search in \texttt{chladni}:\\
tension-curve measurements, solver requirements,\\
and a triage of candidate methods}
\author{Technical report, compiled from \file{notes/} and \file{docs/design/}}
\date{1 September 2026}

\begin{document}
\maketitle

\begin{abstract}
\noindent
This report consolidates a measurement session on the eigenvalue-search layer of
\file{chladni}, the method-of-particular-solutions (MPS) solver for Laplacian eigenvalues of
planar domains. It has three parts. Part~1 reports the first tension-curve measurements taken
on \emph{real} corner domains with corner-adapted Fourier--Bessel bases, rather than on the
surrogate or on a single test shape: a seven-domain roster with bases validated against
tabulated eigenvalues, the finding that the tension plateau is stable across domains
($0.12$--$0.36$) while the attainable floor spreads four to five orders \emph{within} a single
scan, an end-to-end recall-versus-stride measurement showing that the project's standing
``six samples per gap'' rule is wrong by a factor of three, and the identification of the
minimum \emph{local} gap --- $0.13$ to $0.79$ of Weyl's mean --- as the quantity that actually
governs recall. It also records that the Weyl completeness audit reads as a false alarm on
every symmetric domain if it is run before multiplicity estimation.
Part~2 translates those measurements into a requirements list for the minimization and
multiplicity-estimation routines, organized as functional requirements, required diagnostics,
and anti-features that were measured \emph{not} to work, each tied to the measurement that
forces it. It also records that the harness's search interface must widen from a tension
function to a mode function, because the collapse-versus-merge decision needs eigenfunctions.
Part~3 reframes the problem as a level-set/certification problem rather than global
optimization, enumerates 27 candidate methods, and reports six falsifying tests: Weyl-stratified
sampling and the argument principle on $\det A_B$ are dead; pilot-calibrated Lipschitz
certification is \emph{unsound}, not merely loose; the exact $k$-derivative identity holds once
a $k$-dependent column scale is corrected for; DIRECT is strong; and the eigenvector-swing
proposal was falsified and reinterpreted into a cell decomposition. The composed candidate
(swing cells $+$ DIRECT $+$ refinement) reaches mean recall $0.975$ on a five-domain roster at
$2.2\times$ the cost of a uniform grid. A prominent self-correction is recorded: an earlier
claim that the swing separates pairs $290\times$ tighter than the dip was a statement about
\emph{detection}, and had to be repaired with a bisection-toward-the-peak step before it
became a statement about \emph{separation}.
\end{abstract}

\tableofcontents

\newpage

\section{Scope, sources, and what is new}

The layer under discussion sits above \file{mps}. Its contract is
\begin{quote}
given a tension function returning the ascending vector
$\sione \le \sitwo \le \cdots$, find every approximate eigenvalue in $[a,b]$ with its
multiplicity, using as few evaluations as possible,
\end{quote}
and it sees only the tension function and the interval: no domain, no basis, no pencil
\sm{preamble}. That restriction is deliberate --- it is what makes two candidate searches
comparable, and what would let a later finite-element or boundary-integral method reuse the
same search. Part~2, \S\ref{sec:interface}, records the one measured reason to widen it.

The material reported here is drawn from
\file{notes/landscape\_findings.md} \S\S20--25 (new),
\file{notes/solver\_features.md},
\file{notes/approaches.md}, and
\file{docs/design/search.md} \S2. Sections 1--19 of \file{landscape\_findings.md} are earlier
context; they are drawn on for background but not re-reported. The supporting scripts are
\file{notes/curve\_lab.py} (basis builder, scanner, roster),
\file{notes/curve\_study.py} (the three studies of \S20),
\file{notes/adaptive\_scan.py} (\S21),
\file{notes/test\_approaches.py} (\S22),
\file{notes/swing\_pairs.py} (\S23), and
\file{tests/searches.py::swing\_direct\_search} (\S\S24--25).

\subsection{The project-specific machinery, briefly}

Three pieces of vocabulary recur and are worth fixing before the measurements.

\paragraph{Tension and the GSVD pencil.}
MPS assembles, at a trial value $\lam$, a matrix of particular solutions of the Helmholtz
equation evaluated at boundary collocation points ($A_B$) and at interior points ($A_I$). The
\emph{tension} is the generalized singular value spectrum of the pencil $(A_B, A_I)$: it
measures how small the boundary residual of a candidate solution can be made while its
interior norm is held away from zero. Writing $\sione(\lam) \le \sitwo(\lam) \le \cdots$ for
that ascending spectrum, the eigenvalues of the domain appear as local minima of $\sione$. The
GSVD formulation is used in preference to a determinant or a plain least-squares residual
because it is what makes the objective numerically stable; this matters directly to A11 in
Part~3. The interior points are a \emph{normalization}, not a quadrature: their only job is to
keep $\|A_I c\|$ away from zero \sm{1}.

\paragraph{The cone model.}
Near an eigenvalue $\lam_i$ the measured shape of the first tension is
\begin{equation}
  \sione(\lam) \;\approx\; \sqrt{\varepsilon^2 + c^2 (\lam-\lam_i)^2},
  \label{eq:cone}
\end{equation}
a hyperbola that looks like a cone with a rounded bottom. It fits the far field to a median
relative residual of 1--3\% on rect, disk, eq\_tri and the L-shape, failing only where the
floor approaches the plateau, i.e.\ where the basis cannot represent the eigenfunction
\sm{1}. The two parameters are independent and it matters which is which:
$\varepsilon$ is the \emph{floor}, $\sione$ at the eigenvalue, and is a property of
\emph{the basis}; $c$ is the \emph{wall slope}, $d\sione/dk$ away from the eigenvalue, and is a
property of \emph{the domain}. On the L-shape at 20, 40 and 60 Fourier--Bessel terms,
$\varepsilon$ falls $4.0\times10^{-5} \to 1.6\times10^{-8} \to 1.1\times10^{-12}$ while $c$
moves only $1.602 \to 1.595$. The ratio $\varepsilon/c$ is the \emph{resolution limit}: the
scale in $k$ at which the cone flattens into its quadratic bottom, so two eigenvalues closer
than that have overlapping bottoms and cannot be separated. Since $c$ is fixed by geometry,
$\varepsilon$ alone sets it --- which is the concrete form of ``the search is bounded by the
basis''. Away from all eigenvalues $\sione$ sits on a \emph{plateau}, denoted $S$ below; the
plateau is what $\tau_{\mathrm{scan}}$ is measured against.

\paragraph{Multiplicity by (C1) and (C2).}
Multiplicity $\mu$ at a located level is read off the tension spectrum by two conditions:
(C1) a \emph{ratio gap} --- $\sig_{j+1}/\sig_j$ is $1.0$--$1.7$ within a multiplet and
$10^{11}$--$10^{15}$ across the gap to the next branch, so a gap parameter $G=10$ separates
them; and (C2) a \emph{local-minimum} condition --- $\sig_j$ must itself have a local minimum
at that $\lam$, with prominence at least $P$, on a bracket wide enough to see it. (C2) is
load-bearing rather than belt-and-braces: once the separation $d \lesssim \varepsilon/c$, the
ratio rule alone reports multiplicity $2$ at a \emph{simple} eigenvalue, and (C2) rejects it
every time because $\sitwo$ there is a monotone wall \sm{1}, \sfe{F9}.

\section{The governing principle: what a successful run means}
\label{sec:opdef}

Every approach below is judged against \file{docs/design/search.md} \S2, the \emph{operational
definition}. It is summarized here because nothing later is meaningful without it.

The layer above sets \verb|prec|. The contract is that the user gets eigenvalues to relative
accuracy \verb|prec|, or an honest statement that MPS could not reach it. Doing better than
asked is always allowed; doing materially more work than \verb|prec| requires is not.

\subsection{The rule that decides which constants may be fixed}

A constant may be \textbf{fixed} only if it is dimensionless and separates two populations
whose separation does not shrink as \verb|prec| tightens. Anything carrying units of $\lam$,
and anything expressing \emph{how accurate}, must derive from \verb|prec| \sm{2.1}. This rule,
rather than any individual number, is what to apply when a new parameter appears --- and it is
the reason the roster measurement of the plateau in Part~1 is decisive for
$\tau_{\mathrm{scan}}$.

\subsection{Assumptions and ground truth}

The pencil --- basis, collocation points, regularization \emph{and} rank policy --- is given
and adequate. The solver may \emph{detect} inadequacy; it is not responsible for repairing it.
Adequacy is quantitative: $\text{gap} > 20\,\varepsilon/c$ on $[a,b]$, from a direct
measurement in which (C1) is correct at $d/(\varepsilon/c) = 23$ and wrong at $0.45$ and
$0.15$ \sm{2.2}.

The ground truth $S(\verb|prec|)$ is the ordered multiset of local minima of $\sione$, each
repeated by its multiplicity $\mu_i$ (the largest $\mu$ such that $\sig_j$ has a local minimum
at $\lam_i$ for every $j \le \mu$), with two floating-point scales attached: ``has a local
minimum'' means prominence at least $P = 1 + m\eta$ with a floor $P \ge 2$, and ``at $\lam_i$''
means two branch minima lie within $\verb|tol_merge| = \verb|prec|$, relative. $S$ is indexed
by \verb|prec| deliberately: two levels $4\times10^{-9}$ apart are \emph{one} level at
$\verb|prec| = 10^{-8}$ and \emph{two} at $10^{-12}$ \sm{2.3}.

\subsection{The three clauses}

A solver is \textbf{successful at \texttt{prec}} on $[a,b]$ when all three hold \sm{2.4}.

\begin{description}[leftmargin=2.2em,style=nextline]
\item[(1) Completeness and multiplicity.] $\hat S$ matches $S(\verb|prec|)$ as a multiset.
  Among roots reported close together, collapse-versus-merge is decided by the \emph{interior
  Gram angle}, not by distance: with $u_i = I(\hat\lam_i)\,\mathrm{coef}(\hat\lam_i,1)$,
  $|\langle u_1,u_2\rangle| \ge \tau_{\mathrm{gram}}$ means the same eigenfunction (collapse and
  recount) and $< \tau_{\mathrm{gram}}$ means independent (merge and sum). Distance survives
  only as a prefilter deciding which pairs are worth testing.
\item[(2) Accuracy.] Each $|\hat\lam_i - \lam_i|/\lam_i \le \verb|prec|$, \emph{or} the solver
  reports that $\delta_{\mathrm{res}}(\lam_i) > \verb|prec|\cdot\lam_i$ and returns
  $\delta_{\mathrm{res}}$ as the attained bound. There is no predicted-accuracy constant:
  $\delta_{\mathrm{res}}$ is measured from the parabola fit the refinement phase already
  performs.
\item[(3) Honesty.] The reported per-eigenvalue bar contains the true error, and the
  completeness diagnostics fire whenever (1) fails.
\end{description}

Clause~(3) is not decoration. $S$ is unknown on every domain that matters --- \file{L\_shape},
\file{GWW1} and \file{mushroom} have tabulated prefixes only --- so (1) and (2) need a
truth-free proxy on a real run. Part~1 is largely about which proxies actually work.

\subsection{The constants, and not more work than asked}

Fixed (dimensionless, margin grows as \verb|prec| tightens): $P = 1+m\eta$ with floor 2;
$G = 10$; $\tau_{\mathrm{gram}} = 0.5$; the contract $20\,\varepsilon/c$; the Weyl tolerance
$0.5$; and $\tau_{\mathrm{scan}}$ \emph{as a fraction of the plateau, never an absolute
tension}. Derived from \verb|prec|: $\verb|tol_merge| = \verb|prec|$, the accuracy target
$\verb|prec|\cdot\lam$, and the refinement stop
$\max(\verb|prec|\cdot\lam,\ \delta_{\mathrm{res}})$ \sm{2.5}.

Section 2.6 adds that \textbf{completeness is not a precision question} and the scan stride
must not tighten with \verb|prec|. That clause is right; the \emph{number} attached to it in
the same paragraph --- ``about six samples per eigenvalue gap'' --- is what Part~1 falsifies.

Finally, \S2.7 records a real ceiling: $\verb|tol_merge| = \verb|prec|$ against pairs
resolvable to about $3\times10^{-6}$ relative leaves a factor of 30 at $\verb|prec|=10^{-8}$,
but below $\verb|prec|\sim 4\times10^{-9}$ it \emph{inverts} --- the discretization's own
splitting of a degenerate pair ($1.8\times10^{-9}\lam$ for plane waves,
$3.9\times10^{-9}\lam$ for Fourier--Bessel) exceeds $\verb|tol_merge|$ and a true double is
reported as two simple eigenvalues.

\part*{Part 1: Tension curves on real corner domains}
\addcontentsline{toc}{section}{Part 1: Tension curves on real corner domains}
\setcounter{section}{2}

\section{The roster and how the bases were built and validated}
\label{sec:roster}

Everything measured before \S20 used either the surrogate, a rectangle with plane waves, or the
L-shape. The solver has to work on domains with \emph{several} reentrant corners, where
Fourier--Bessel families at those corners are what make it work at all \lf{20}.

The first thing \file{curve\_lab.py} had to get right is which corners need such a family. A
corner is \textbf{singular} when its local orders $(\ell+\text{offset})\pi/\omega$ are not all
integers, which under a single boundary condition is just ``$\pi/\omega$ is not an integer''.
\verb|Corner.is_singular| then meant ``geometric corner'', a much larger and different set ---
all four corners of a rectangle were \verb|is_singular| and none of them needs a
Fourier--Bessel family, so the basis builder keyed on the fractional exponent itself \lf{20}.
That property is gone. The question is now asked where its answer actually lives, as
\verb|CornerModes.is_singular| (the offset comes from the boundary-condition pair through
Grisvard's Robin rule, which is not a geometric fact); \verb|geometry| keeps only the angle
half of it, \verb|Corner.is_reflection_closed|, meaning $\omega=\pi/m$. The builder's test is
unchanged and is now one call: \verb|CornerModes.from_corner(c).is_singular|.

Bases were built as plane waves sized by \verb|n_useful(k*R)| plus 40 Fourier--Bessel terms per
non-analytic corner. Crucially, each basis was \emph{validated before} anything was concluded
from its curves, by evaluating the tension at the tabulated eigenvalues
(Table~\ref{tab:roster}).

\begin{table}[htbp]
\centering
\caption{The roster. Columns are basis size, corner structure, and the tension attained at the
tabulated eigenvalues --- the basis-adequacy check run before any curve was interpreted
\lf{20}.}
\label{tab:roster}
\begin{tabular}{lrll}
\toprule
domain & cols & corner structure & tension at table values \\
\midrule
\file{L\_shape}  & 134 & 1 reentrant                        & $10^{-8}\,..\,2\times10^{-11}$ \\
\file{mushroom}  & 168 & 2 reentrant                        & $7\times10^{-10}\,..\,2\times10^{-10}$ \\
\file{GWW1}      & 272 & 2 reentrant $+$ two $3\pi/4$       & $3\times10^{-9}\,..\,10^{-9}$ \\
\file{H\_shape}  & 258 & 4 reentrant                        & $5\times10^{-8}\,..\,2\times10^{-8}$ \\
\file{plus}      & 248 & 4 reentrant, $D_4$                 & (no table) \\
\file{chevron}   & 252 & an $18^\circ$ spike                & $2\times10^{-7}\,..\,5\times10^{-8}$ \\
\file{rect81}    & 176 & no singular corner, control        & (closed form) \\
\bottomrule
\end{tabular}
\end{table}

\file{cut\_square} was dropped from the roster. The \verb|reference| recipe for it is an
\emph{origin-centred} Fourier--Bessel expansion of 320 terms, which the corner-adapted builder
does not offer, and plane waves leave the tension at $10^{-2}$. That is a basis-planning gap,
not a curve finding \lf{20}, and it is the reason the roster contains no curved boundary at all
(see \S\ref{sec:open}).

\section{The plateau is stable across domains; the floor is not, within one scan}
\label{sec:plateau}

\begin{table}[htbp]
\centering
\caption{Plateau height and the range of attained floors within a single scan \lf{20.1}.}
\label{tab:plateau}
\begin{tabular}{lrrrr}
\toprule
domain & plateau & floor min & floor max & spread \\
\midrule
\file{L\_shape}  & $0.1981$ & $3.6\times10^{-11}$ & $9.1\times10^{-7}$ & $2.5\times10^{4}$ \\
\file{mushroom}  & $0.1829$ & $2.2\times10^{-10}$ & $1.4\times10^{-8}$ & $6.4\times10^{1}$ \\
\file{GWW1}      & $0.1228$ & $1.0\times10^{-10}$ & $2.6\times10^{-8}$ & $2.5\times10^{2}$ \\
\file{H\_shape}  & $0.2251$ & $2.7\times10^{-12}$ & $1.2\times10^{-7}$ & $4.4\times10^{4}$ \\
\file{plus}      & $0.2131$ & $1.1\times10^{-11}$ & $3.7\times10^{-8}$ & $3.4\times10^{3}$ \\
\file{chevron}   & $0.3555$ & $4.9\times10^{-8}$  & $2.3\times10^{-7}$ & $4.7\times10^{0}$ \\
\file{rect81}    & $0.1541$ & $5.3\times10^{-7}$  & $4.9\times10^{-2}$ & $9.2\times10^{4}$ \\
\bottomrule
\end{tabular}
\end{table}

Two findings, and they point in opposite directions.

\paragraph{The plateau spans 0.12 to 0.36 across seven domains} --- a factor of $2.9$, across
areas from $1.0$ to $14.0$ and from zero to four reentrant corners. That is the strongest
support yet for $\tau_{\mathrm{scan}}$ as a \emph{fraction of the plateau}, and against any
absolute tension threshold \lf{20.1}. Under the \S\ref{sec:opdef} rule, the plateau-relative
form is admissible as a fixed constant and the absolute form is not.

\paragraph{The floor spreads four to five orders \emph{within} a single scan on three of seven
domains.} The L-shape's five orders, reported earlier \sm{1} and explained there by modes whose
eigenfunction vanishes at the reentrant corner and is therefore analytic there, was not
special. This is the direct justification for clause~(2) of the operational definition being
\emph{per eigenvalue}: one \verb|prec| target cannot be met uniformly when the attainable floor
moves by $10^{5}$ across the interval being scanned \lf{20.1}.

\file{rect81} is the informative failure. Its floor maximum is $4.9\times10^{-2}$, which is no
dip at all --- about a third of its own plateau. A high aspect ratio with plane waves is a
genuinely inadequate basis, and \file{rect81} is the only roster member whose audit flags for
that reason rather than for multiplicity \lf{20.1}. It recurs throughout Part~1 as the domain
that no amount of sampling can fix.

\section{Recall versus stride: six samples per gap is not enough}
\label{sec:stride}

\file{search.md} \S6 states that ``roughly six samples per eigenvalue gap always place three
inside every dip, and Weyl's law sizes that directly''. That was \emph{inferred} from the
width/gap invariant on four benign domains. It was then measured end to end: subsample the fine
scan at stride $h$ and count how many true dips survive as discrete minima
(Table~\ref{tab:stride}).

\begin{table}[htbp]
\centering
\caption{Recall as a function of stride, measured by subsampling the fine scans \lf{20.2}.
The first column is the coarsest stride, in samples per median gap, at which recall is still
exactly $1$.}
\label{tab:stride}
\begin{tabular}{lrrr}
\toprule
domain & recall $=1.000$ down to & at $\sim6$ per gap & at $\sim4$ per gap \\
\midrule
\file{L\_shape}  & $4.8$ per gap & $1.000$ & $0.889$ \\
\file{mushroom}  & $5.0$         & $1.000$ & $0.917$ \\
\file{GWW1}      & $8.0$         & $0.933$ & $0.867$ \\
\file{H\_shape}  & $9.1$         & $0.875$ & $1.000$ \\
\file{plus}      & $4.4$         & $1.000$ & $1.000$ \\
\file{chevron}   & $10.1$        & $0.667$ & $0.667$ \\
\file{rect81}    & $17.7$        & $0.929$ & $0.929$ \\
\bottomrule
\end{tabular}
\end{table}

\textbf{Full recall on every domain needs about 18 samples per median gap, not six. At six,
four of seven domains lose eigenvalues} \lf{20.2}. This is a direct falsification of a standing
design rule, and it is the single most consequential number in Part~1: the rule it replaces is
cited in \file{search.md} \S2.6 as part of the ``not more work than asked'' argument, and the
replacement is three times as expensive.

Two qualifications are recorded in the source and are preserved here. First, 18 is a
\emph{fitted} number on seven domains, not a bound \sfe{6}. Second, the $\sim4$-per-gap column
is not monotone with the $\sim6$ column on \file{H\_shape} ($1.000$ against $0.875$), which is
a reminder that at coarse strides recall is a sampling accident rather than a smooth function
of the stride.

\subsection{The mechanism: minimum local gap, not Weyl's mean}

The width/gap invariant is about the \emph{mean} gap; recall is bounded by the \emph{minimum
local} gap. Measured, $\min(\text{observed gap})/\text{Weyl mean gap}$:

\begin{table}[htbp]
\centering
\caption{Minimum observed gap as a fraction of the Weyl mean gap \lf{20.2}.}
\label{tab:mingap}
\begin{tabular}{lrlr}
\toprule
domain & ratio & domain & ratio \\
\midrule
\file{mushroom} & $0.127$ & \file{chevron}  & $0.336$ \\
\file{H\_shape} & $0.210$ & \file{L\_shape} & $0.571$ \\
\file{rect81}   & $0.295$ & \file{plus}     & $0.740$ \\
                &         & \file{GWW1}     & $0.792$ \\
\bottomrule
\end{tabular}
\end{table}

\textbf{The local gap runs $1.3\times$ to $8\times$ tighter than Weyl's mean.} A stride computed
from Weyl is therefore an estimate of the wrong quantity, and no a-priori stride can be safe
--- which is the case for correcting the stride \emph{in the loop} rather than choosing it once
\lf{20.2}. This finding does more work than the 18-per-gap number itself: it is the reason
Part~3's A2 (Weyl-stratified sampling) was tested at all, and the reason its falsification kills
the whole a-priori-schedule family.

A third measurement in the same section corrects an earlier invariant. The half-width/gap ratio
is \emph{not} the near-constant reported in \S14 of the notes. Measured consistently as the
half-width at half the plateau, it runs $0.13$ to $0.39$ across the roster --- a $3\times$
range, with \file{H\_shape} and \file{chevron} at the top. Section~14's $0.46$--$0.52$ used a
different definition (full width at the plateau) on four benign domains; converted, the roster
runs about $0.5$ to $1.6$. \textbf{The invariant is much looser than one domain family
suggested} \lf{20.2}.

\section{The Weyl audit is a post-multiplicity check}
\label{sec:weyl}

The Weyl deficit is the only completeness diagnostic available on a domain with no closed form:
compare the number of levels found below a cut against the three-term asymptotic count
$N(\lam)$. When the audit was run at each stride, it flagged \verb|MISS| at recall $1.000$ on
four of seven domains. That looks like a false positive and is not:

\begin{quote}
$\sione$ shows one dip per \textbf{level}, and Weyl counts \textbf{with multiplicity}. The
deficit was reporting exactly the multiplicity that had not been estimated yet \lf{20.3}.
\end{quote}

Closing the loop --- (C1) with $G=10$ plus (C2) at prominence $P=2$ on a wide bracket, then
auditing the \emph{unfolded} list --- gives Table~\ref{tab:weyl}.

\begin{table}[htbp]
\centering
\caption{Weyl deficit before and after multiplicity estimation \lf{20.3}.}
\label{tab:weyl}
\begin{tabular}{lrrrr}
\toprule
domain & dips & $\sum \mu$ & deficit on dips & deficit after multiplicity \\
\midrule
\file{L\_shape} &  9 & 10 & $0.41$ & $0.41$ \\
\file{mushroom} & 12 & 12 & $0.59$ & $0.59$ \\
\file{GWW1}     & 15 & 15 & $0.33$ & $0.33$ \\
\file{H\_shape} &  8 &  8 & $1.18$ & $1.18$ \\
\file{plus}     &  8 & 10 & $2.79$ & $\mathbf{0.79}$ \\
\file{chevron}  &  3 &  3 & $0.12$ & $0.12$ \\
\file{rect81}   & 14 & 14 & $2.55$ & $2.55$ \\
\bottomrule
\end{tabular}
\end{table}

\file{plus} is the case that proves it: $D_4$ symmetry, two doubles found, and the deficit falls
from $2.79$ to $0.79$. \textbf{(C1)$+$(C2) works on a domain with four reentrant corners and a
corner-adapted basis} --- the earlier 9-for-9 multiplicity score was on the disk and the
equilateral triangle, both smooth \lf{20.3}.

The residual flags are then real, and each names a different fault:
\begin{itemize}[nosep]
\item \file{rect81} at $2.55$ --- an inadequate \emph{basis}, floor $4.9\times10^{-2}$ at some
  levels;
\item \file{H\_shape} at $1.18$ --- a genuine miss; its minimum observed gap is $0.378$ against a
  Weyl mean of $1.795$, so there is a tight pair the stride merged;
\item \file{plus} at $0.79$ --- one more level or multiplicity not yet accounted for.
\end{itemize}

And \file{chevron} reads $0.12$ at recall $0.667$. \textbf{The audit is blind here}, because
three distinct levels give too few cuts and every cut lies below the top level. That limitation
is documented in \verb|asymp.weyl_deficit|, and this is it firing on a real domain \lf{20.3}.

\subsection{What \S20 changes}

\begin{itemize}[nosep]
\item $\tau_{\mathrm{scan}}$ as a fraction of the plateau is now measured across seven domains,
  not assumed.
\item The per-eigenvalue accuracy clause is necessary, not defensive: $10^5$ floor spread in one
  scan.
\item \textbf{The ``six samples per gap'' rule in \file{search.md} \S6 is wrong and should be
  replaced.} About 18 per median gap reaches full recall on this roster, and even that is a
  fitted number rather than a bound, because it is the minimum local gap that governs.
\item \textbf{The Weyl audit must run on the unfolded list, after multiplicity.} Run on dips, it
  reads as a false alarm on every symmetric domain --- which is exactly the class the solver will
  meet at the optimizers of a shape optimization.
\item The audit needs enough distinct levels to have cuts. Below about five it says nothing, and
  \file{chevron} shows it saying nothing while a third of the spectrum is missing \lf{20.4}.
\end{itemize}

\section{The Weyl deficit as a localizer, not a verdict}
\label{sec:localizer}

Since no a-priori stride is safe, the stride has to be corrected in the loop.
\file{notes/adaptive\_scan.py} tests one way: make the Weyl deficit a \emph{profile} over cuts,
so it says not only \emph{that} a level is missing but \emph{where}, and re-scan only that
interval. Cost then goes as the number of tight spots rather than as the whole range \lf{21}.

Scored against the \S20 fine scans as truth, with $s_0 = 8$ (one sample per Weyl gap divided by
8), prominence $P=2$, $G=10$, $\tau_{\mathrm{scan}} = 0.5\,S$:

\begin{table}[htbp]
\centering
\caption{The audit-guided adaptive scan \lf{21}. ``found'' includes spurious roots; the fine
scan's evaluation count is given for comparison.}
\label{tab:adaptive}
\begin{tabular}{lrrrrrrr}
\toprule
domain & truth & found & recall & spurious & deficit & evals & fine-scan $n$ \\
\midrule
\file{L\_shape} &  9 &  9 & $1.000$ & 0 & $0.13$ &  205 & 317 \\
\file{mushroom} & 12 & 11 & $0.917$ & 0 & $0.47$ &  238 & 348 \\
\file{GWW1}     & 15 & 15 & $0.867$ & 2 & $-0.09$&  326 & 483 \\
\file{H\_shape} &  8 &  8 & $0.875$ & 1 & $0.81$ & 1556 & 376 \\
\file{plus}     &  8 &  8 & $0.875$ & 1 & $0.43$ &  206 & 355 \\
\file{chevron}  &  3 &  3 & $0.667$ & 1 & $-0.78$&   84 & 154 \\
\file{rect81}   & 14 & 17 & $1.000$ & 0 & $0.94$ & 4094 & 582 \\
\bottomrule
\end{tabular}
\end{table}

\subsection{The bug it found first, which was the author's}

The first version anchored the profile by subtracting $N(a)$, on the reasoning that the count is
only meaningful relative to the bottom of the scanned range. \textbf{The three-term count is an
asymptotic and goes negative at the low end} --- measured $N(a)$ of $-0.230$ on
\file{L\_shape}, $-0.687$ on \file{H\_shape}, $-0.752$ on \file{rect81}, $-1.041$ on
\file{chevron}. Subtracting it \emph{adds} that much to every deficit, so the loop chased a
phantom:

\begin{table}[htbp]
\centering
\caption{Effect of removing the spurious anchor \lf{21}.}
\begin{tabular}{l@{\quad}p{9.2cm}}
\toprule
\file{plus}     & 1794 evaluations $\to$ 206, an $8.7\times$ reduction, and the deficit
                  reads $0.43$ instead of $0.89$ \\
\file{L\_shape} & deficit $0.36 \to 0.13$ \\
\file{GWW1}     & deficit $-0.01 \to -0.09$ \\
\bottomrule
\end{tabular}
\end{table}

The count is absolute, so the correct form subtracts nothing and instead \emph{requires} the
found set to be a complete prefix from the bottom of the spectrum --- the same precondition
\verb|asymp.weyl_deficit| already carries, and it should be asserted rather than papered over
with an anchor that looks principled and is not \lf{21}.

\subsection{Two structural flaws, both real}

\paragraph{(1) A cluster miss is \emph{at} a level, not \emph{between} levels.}
The localizer re-scans the interval between the two found levels straddling a hot cut. But when
the missing eigenvalue is the second member of an unresolved pair, it sits at the \emph{same}
location as a found level, and no amount of re-scanning between levels will ever reach it. That
is \file{mushroom} (recall $0.917$, deficit $0.47$, minimum observed gap $0.352$ against a Weyl
mean of $2.771$) and \file{H\_shape} ($0.875$, deficit $0.81$, minimum gap $0.378$ against
$1.795$). \file{H\_shape} became \emph{more} expensive after the anchor fix, $958 \to 1556$,
because the corrected deficit is real and the loop keeps trying to fix it by the one mechanism
that cannot. The deficit response therefore has to be two-pronged:
\begin{center}
\begin{tabular}{ll}
\toprule
between levels & re-scan the interval \quad(a genuinely unsampled dip) \\
at a level     & re-examine multiplicity \quad(an unresolved cluster) \\
\bottomrule
\end{tabular}
\end{center}
and only the first is implemented \lf{21}.

\paragraph{(2) No basis-adequacy termination guard.}
\file{rect81} spends 4094 evaluations --- $7\times$ the fine scan --- and still reports a
deficit of $0.94$, because its floor reaches $4.9\times10^{-2}$ at some levels and no stride can
fix a basis. This is the unbounded-refinement stall, and the predicate that stops it already
exists: $\text{gap} > 20\,\varepsilon/c$. Measured on \file{H\_shape} at every level,
$\text{gap}/(\varepsilon/c)$ runs $2.2\times10^{6}$ to $1.4\times10^{11}$ --- massively adequate,
so the guard correctly does \emph{not} fire there and correctly separates ``refine harder'' from
``the basis cannot do this''. That discrimination is the thing the loop most needs and it is one
predicate away \lf{21}.

\subsection{Two cautions on the scoring itself}

\file{chevron} reads deficit $-0.78$ at recall $0.667$: the audit says clean while a third of the
spectrum is missing. Inside a control loop this is worse than in \S\ref{sec:weyl}, because
\textbf{the loop terminates on it}.

And the truth is basis-dependent. ``Truth'' here is the fine scan's dips below $0.5\,S$. On
\file{rect81} the floor reaches $4.9\times10^{-2}$, about a third of the plateau, so real dips
fall below that cutoff and the adaptive scan finding 17 against a ``truth'' of 14 is the truth
undercounting, not the search over-reporting. \textbf{Any recall number on a domain whose basis
is marginal is measuring the cutoff as much as the search} \lf{21}.

\paragraph{Verdict.} The localizer is worth keeping --- it is what turned \file{plus} into 206
evaluations, and on the three domains where the miss is genuinely a resolution problem it reaches
the fine scan's recall at $0.6\times$ its cost. It is not sufficient on its own; the two additions
it needs (the at-level prong, and the $20\,\varepsilon/c$ guard) are both specified and neither
is speculative. What it should \emph{not} be is a verdict \lf{21}.

\part*{Part 2: What the solver has to have}
\addcontentsline{toc}{section}{Part 2: What the solver has to have}

\section{Functional requirements}
\label{sec:functional}

The list below is from \file{notes/solver\_features.md}, derived from the operational definition
of \S\ref{sec:opdef} and from the measurements of Part~1. Every entry names the measurement that
forces it.

\begin{description}[leftmargin=1.2em,style=nextline,itemsep=0.4em]

\item[F1. Recall first, and recall is conditional.]
Missing an eigenvalue cannot be repaired downstream; a located root a digit short can. Every
claim of completeness is conditional on the scan's own assumptions and must say so.

\item[F2. Adaptive subdivision, not a uniform grid.]
A pair at separation $d$ is resolved once sampling near it is finer than about $d/3$. A uniform
grid pays that everywhere at $n \sim 3(b-a)/d$; subdivision pays $O(\log(1/d))$. Measured
$2048$ against $47$ evaluations at $d = 3\times10^{-3}$, and $6\times10^{10}$ against $97$ at
$d = 10^{-10}$. $d$ is not known in advance, so the uniform cost is paid against a pair that may
not exist. \sm{3}

\item[F3. A stride that is corrected in the loop.]
Six samples per gap loses eigenvalues on four of seven real domains; full recall needs about 18
per median gap. And the governing quantity is the \emph{minimum local} gap, which runs $0.13$ to
$0.79$ of Weyl's mean --- so no a-priori stride is safe and the correction has to happen during
the solve. \lf{20.2}, and \S\ref{sec:stride} above.

\item[F4. Per-solve calibration of $c$.]
$c$ runs $0.45$ to $1.73$ across domains, stable within one. A fixed Lipschitz constant is either
wrong or wasteful. \lf{2}

\item[F5. Threshold rank, re-derived at each $\lam$.]
The numerical rank \emph{grows} with $\lam$, so a rank frozen at one point is unsafe at the low
end of a range: 39 columns chosen at $k = 3.95$ gives 85 spurious minima at $k = 2.35$, where the
correct 31 gives exactly one. Threshold cannot pick an unsafe rank by construction, at
$1.2$--$1.6\times$ the cost. \lf{3, 5, 8}

\item[F6. Two-level rank from \emph{one} factorization.]
The scan wants a low rank (smooth, wide dips); the refinement wants a high one (accurate floor).
The L-shape locates $\lam_1$ to $3.2\times10^{-12}$ at rank 111 against $1.4\times10^{-10}$ at
rank 61, because what limits it is the floor rather than the roughness. One SVD serves both ---
no second basis, no second assembly. \lf{6} \textbf{Not built.}

\item[F7. Refine in $\lam$.]
Fitting $\sione^2$ over identical physical points, the residual is smaller in $\lam$ in 17 of 17
cases across four domains, median $5.1\times$, and the located vertex is $5.1\times$ more
accurate. The \emph{scan} variable is a separate choice. \lf{3b}

\item[F8. Squared objective.]
Minimize $\sione^2$, never $\sione$: squaring removes the kink and makes refinement
machine-exact. \sm{3}

\item[F9. Multiplicity by (C1) \emph{and} (C2).]
The ratio gap alone double-counts once $d \lesssim \varepsilon/c$, reporting multiplicity 2 at a
\emph{simple} eigenvalue; (C2) rejects it every time because $\sitwo$ there is a monotone wall.
Confirmed on a domain with four reentrant corners: on \file{plus\_shape} the two $D_4$ doubles are
found and the Weyl deficit falls $2.79 \to 0.79$. \lf{9.4, 20.3}

\item[F10. Collapse-versus-merge by the interior Gram angle.]
Two roots at the same separation can be one eigenvalue found twice or two genuine ones; no length
threshold decides it. Distinct pairs read $|\cos| \le 0.10$, duplicates $\ge 0.948$. Within one
$\lam$ the test is \emph{vacuous} --- the GSVD makes the interior images orthogonal by
construction, measured $4.6\times10^{-16}$ --- so it is only informative \emph{across} roots.
\lf{15}

\item[F11. Work scaled to \texttt{prec}.]
Refinement stops at $\max(\verb|prec|\cdot\lam,\ \delta_{\mathrm{res}})$; the first term binds
whenever the basis is better than the request. Completeness is \emph{not} a precision question
and the stride must not tighten with \verb|prec|. \sm{2.6}

\item[F12. Deterministic output.]
The same domain returns the same eigenvalue twice. An unseeded interior fill gives a $110\times$
run-to-run spread in the located eigenvalue and puts that into the shape objective as noise.
\lf{17}

\end{description}

\section{Required diagnostics}
\label{sec:diagnostics}

These are output, not instrumentation. $S$ is unknown on every domain that matters, so they are
the only signals that F1 has failed.

\begin{description}[leftmargin=1.2em,style=nextline,itemsep=0.4em]

\item[D1. The Weyl deficit, on the \emph{unfolded} list.]
The only completeness check that does not depend on the search being right. It found the
recursion-floor bug that prefix truth did not. Two hard limits, both measured firing on real
domains: run \emph{before} multiplicity it reads as a false alarm on every symmetric domain
($2.79$ on \file{plus}, which is really $0.79$), and below about five distinct levels it says
nothing at all while a third of the spectrum is missing (\file{chevron}, $0.12$ at recall
$0.667$). \lf{20.3}

\item[D2. $L$ and the largest observed secant slope, with their ratio.]
A small margin means the constant was being pushed and the run deserves a rerun. The only signal
that the first pass was too coarse. (See A6 in Part~3: this diagnostic is what would have caught
the unsound pilot.)

\item[D3. Per root: $\delta_{\mathrm{res}}$, and which termination bound was active.]
This is what makes clause~(2) enforceable while its constant is uncalibrated.

\item[D4. The basis-adequacy predicate, per interval.]
$\text{gap} > 20\,\varepsilon/c$, from a direct measurement --- (C1) correct at
$d/(\varepsilon/c) = 23$, wrong at $0.45$ and $0.15$. \lf{9.4}

\item[D5. A flag when \texttt{tol\_merge} approaches the local resolution limit.]
At $\verb|prec| \sim 4\times10^{-9}$ the discretization's own splitting of a degenerate pair
exceeds \verb|tol_merge| and a true double is reported as two simple eigenvalues. \lf{19}

\item[D6. Evaluation count and attic revivals.]

\end{description}

\section{Anti-features: measured not to work}
\label{sec:anti}

These carry their own numbering in \file{solver\_features.md} (A1--A9) which is \emph{unrelated}
to the approach numbering A1--A27 of Part~3. They are relabelled X1--X9 here to avoid the
collision.

\begin{description}[leftmargin=1.2em,style=nextline,itemsep=0.4em]

\item[X1. No absolute tension threshold, anywhere.]
The floor spreads four to five orders \emph{within} one scan on three of seven domains.
Thresholds are fractions of the plateau, which is $0.12$ to $0.36$ across the whole roster.
\lf{20.1}

\item[X2. No $\sitwo$-\emph{value} gate for cluster detection.]
$\sitwo$'s minimum is $c\,d/2$ only at the crossing point, which is never a sample, so a grid of
spacing $h$ sees $c\,h/2$. The gate tests the \emph{grid}, not the pair. \lf{11}

\item[X3. No offset detector expecting early warning.]
The $\sitwo$ offset relation is real (1\% of $d/2$, confirmed over six decades) but it fires at
exactly the same coarseness at which $\sione$ first shows two minima. It measures $d$ once
resolved; it does not detect earlier. \lf{9.2, 16} (Part~3, \S\ref{sec:swingpairs}, records the
one detector that \emph{does} fire earlier, and the extent to which the corresponding general
claim had been over-stated.)

\item[X4. No distance-only collapse/merge.] See F10.

\item[X5. No structured interior fill.]
\verb|method="grid"| is $30\times$ to $10^{4}\times$ more accurate as a discrete inner product
\emph{and} annihilates the mode $\mathrm{side}+1$ exactly, to $10^{-56}$, because its points sit
on that mode's nodal lines. The interior block wants \emph{non-degeneracy}, not accuracy.
\lf{18}

\item[X6. No interior-cubature project.]
Accurate interior inner products for eigenfunctions are hard and the path is dysfunctional. The
interior block only has to bound $\|u\|$ away from zero. Anything that must be \emph{accurate}
belongs on the Rellich boundary integrals.

\item[X7. No fixed tolerance that expresses resolution or accuracy.] \sm{2.1}

\item[X8. No auto-swap to a symmetry-reduced solve on detected symmetry.]
A tolerance threshold makes the objective discontinuous in shape, and it fires in the endgame of
the optimization. (\file{mps\_symmetry.md} \S2.) This is the constraint that makes A25 in Part~3
a \emph{declared}-symmetry option only.

\item[X9. Do not design against a narrow dip.]
No domain produces one; measured $c$ is $0.445$ to $1.725$ over 19 eigenvalues. The real failure
is a \emph{cluster}. \lf{14}

\end{description}

\section{The interface has to widen}
\label{sec:interface}

The current harness contract is \verb|search(tensions, a, b) -> [(lam, mult)]|, and
\file{search.md} defends it: the search sees only the tension function, which is what makes
candidates comparable and lets a later FEM or boundary-integral method reuse it.

\textbf{F10 breaks that.} The Gram angle needs the candidate eigen\emph{functions}, which
\verb|tensions(x)| cannot supply. Two ways out \sfe{4}:
\begin{enumerate}[nosep]
\item Widen the interface to \verb|modes(x) -> (sigma[:], U[:, :])|, where $U$'s columns are the
  candidate functions sampled on a \emph{fixed} common point set. FEM and BIM produce
  eigenvectors too, so this stays method-agnostic and keeps candidates comparable. The fixed
  point set is already what \verb|MPSPencil| has --- \verb|interior| is set at construction and
  only $k$ changes.
\item Move the collapse phase above the search, into \verb|MPSEigensolver|, as post-processing.
\end{enumerate}
\textbf{(1) is better}, and it is a small change: the harness's \verb|Case| grows one callable,
and \verb|uniform_grid_search| and the lappy baselines ignore it. (2) makes the collapse
decision un-benchmarkable, which is precisely the decision that most needs benchmarking.

This finding is more consequential than it first looks, because Part~3's two surviving approaches
--- the eigenvector swing and the bisection repair --- both read \verb|modes(x)| and neither is
expressible under the narrow contract.

\section{Ranked build order}

\begin{enumerate}[nosep]
\item \textbf{The stride rule and the in-loop Weyl correction (F3, D1).} The measured defect is
  here --- six per gap loses eigenvalues on four of seven domains --- and recall is the headline.
\item \textbf{The interface widening (\S\ref{sec:interface}), then F10.} Everything about
  multiplicity is un-benchmarkable until the harness can see eigenfunctions.
\item \textbf{\texttt{prec}-derived termination (F11)} and the $\delta_{\mathrm{res}}$ report
  (D3). This is what makes the precision contract real rather than aspirational; nothing
  currently reads \verb|prec| at all.
\item \textbf{Two-level rank (F6).} A pure win already measured, and cheap: one factorization.
\item Basis planning, which bounds all of it and is a separate project.
\end{enumerate}

\part*{Part 3: Candidate methods}
\addcontentsline{toc}{section}{Part 3: Candidate methods}

\section{The reframing that generates the list}
\label{sec:reframe}

\textbf{This is not global optimization.} Global optimization finds $\arg\min f$. What is needed
is
\[
  \{\, x \in [a,b] \;:\; x \text{ is a local minimum of } \sione \text{ with prominence} \ge P \,\}
\]
with a multiplicity at each, and a certificate that none was missed. That is a
\textbf{level-set / sublevel-set problem} --- find and certify $\{f \le \tau\}$ --- with a
local-structure problem inside each component \ap{0}. Two consequences follow, and they organize
everything below:
\begin{itemize}[nosep]
\item the hard half is the \textbf{certificate} (what is provably absent), not the location;
\item methods designed to converge fast to \emph{one} minimum are mostly irrelevant, and methods
  designed to \textbf{count} or \textbf{enclose} are directly relevant even when they locate
  nothing.
\end{itemize}
The enumeration is therefore organized by \emph{what supplies the certificate}.

\section{The enumeration}
\label{sec:enum}

\subsection{Family 1 --- deterministic sampling, local extraction}
Certificate: ``the samples were fine enough.'' The weakest kind, and everything currently built
lives here.
\begin{description}[leftmargin=1.2em,style=nextline,itemsep=0.25em]
\item[A1. Uniform grid.] The control, \verb|harness.uniform_grid_search|. Falsifiable claim: a
  grid at $s$ samples per Weyl gap reaches recall $1.0$ for $s \ge s^\ast$. Measured
  $s^\ast \approx 18$ on the roster, and it is not a bound (\S\ref{sec:stride}).
\item[A2. Weyl-stratified grid.] Sample equispaced in $N(\lam)$, not in $\lam$. \textbf{Tested.}
\item[A3. Recursive bisection into every dip.] \verb|lappy_repaired_search|, 10/12 at 5573
  evaluations. The open recursion-floor bug lives here.
\item[A4. Audit-guided local refinement.] \S\ref{sec:localizer}. Works where the miss is a
  resolution problem; structurally cannot reach a miss that sits \emph{at} a level; no basis
  guard.
\item[A5. Coarse-to-fine cascade sharing evaluations.] One sorted cache, repeated passes at
  halved stride, extraction from the cache each time.
\end{description}

\subsection{Family 2 --- bounding and certification}
Certificate: a proven lower bound on $f$ over an interval. The only family that can certify
\emph{absence} rather than assert it.
\begin{description}[leftmargin=1.2em,style=nextline,itemsep=0.25em]
\item[A6. Piyavskii--Shubert with a calibrated $L$.] The originating proposal. Built, broken,
  7/12. \textbf{Tested.}
\item[A7. DIRECT (dividing rectangles).] Lipschitz-free: no constant to calibrate, considers all
  Lipschitz constants at once by keeping the potentially-optimal set. \textbf{Tested.}
\item[A8. Cone-model envelope.] Bound from the measured structure \eqref{eq:cone} rather than from
  secants: between two samples the cone model says how deep a dip can hide.
\item[A9. Interval arithmetic / rigorous enclosure.] The only route to a true proof. Needs an
  interval extension of a generalized singular value through a GSVD --- not available and not
  plausible to build. \textbf{Rejected}, with the reason recorded.
\item[A10. Derivative-bounded envelope.] A two-point Hermite bound, far tighter than a secant
  bound. Gated on A14.
\end{description}

\subsection{Family 3 --- analytic root counting}
Certificate: a contour integral counts the roots inside, exactly and with multiplicity. The only
family that answers the completeness question \emph{directly} rather than by sampling.
\begin{description}[leftmargin=1.2em,style=nextline,itemsep=0.25em]
\item[A11. Argument principle on $\det A_B(\lam)$.] With a square collocation
  ($n_{\mathrm{bdry}} = n_{\mathrm{cols}}$), $\det A_B(\lam)$ is analytic in $\lam$ and
  $(2\pi i)^{-1}\oint d(\log\det)$ counts its zeros. \textbf{Tested.}
\item[A12. Beyn / Sakurai--Sugiura contour integration.] The standard tool for ``all eigenvalues
  inside a contour, with multiplicity'' for nonlinear eigenproblems $T(\lam)v=0$. Obstacle: MPS
  is rectangular and least-squares, so the moment matrices are not defined; a square
  reformulation exists only in A11's setting, which gives up least-squares stability.
\item[A13. Inertia / Sylvester count.] $N(\lam)$ exactly, by counting negative eigenvalues of a
  shifted Hermitian pencil. \textbf{Rejected}: FEM and BIM have this; MPS does not, because there
  is no Hermitian matrix whose inertia counts Laplacian eigenvalues below $\lam$. It is recorded
  because it is the single feature whose absence makes this problem hard, and a hybrid borrowing
  a coarse FEM count purely as a counting oracle is a real option (A24).
\end{description}

\subsection{Family 4 --- derivative-based}
\begin{description}[leftmargin=1.2em,style=nextline,itemsep=0.25em]
\item[A14. Exact $d\sig/d\lam$ from existing primitives.] Every particular solution here is
  $\psi(x) = \Psi(kx)$ about its own centre --- plane waves $e^{ik\,d\cdot x}$ and
  Fourier--Bessel $J_\nu(kr)e^{i\nu t}$ both --- so Euler's identity gives
  $\partial\psi/\partial k = (x\cdot\nabla\psi)/k$ exactly, and $x\cdot\nabla\psi$ is the
  \emph{Rellich multiplier}, which the boundary-integral layer is already planned around.
  \textbf{Tested.}
\item[A15. Sign-of-derivative bracketing.] A bracket is a sign change --- no extra evaluations, and
  no reliance on three-point local-minimum tests that fail in the flat bottom.
\item[A16. Curvature for $\delta_{\mathrm{res}}$.] The second derivative at the minimum gives
  $\delta_{\mathrm{res}}$ directly rather than through a three-point parabola fit.
\end{description}

\subsection{Family 5 --- structure of the branch family}
\begin{description}[leftmargin=1.2em,style=nextline,itemsep=0.25em]
\item[A17. Track all $\sig_j$ as smooth curves.] $\sione$ is $\min_j$ of branches and has kinks at
  crossings; the individual branches are smooth. Continuation turns a non-smooth minimization into
  several smooth ones.
\item[A18. $\sitwo$ as a between-dip certificate.] Between two eigenvalues $\sione$ is the meeting
  of two cones and $\sitwo$ is bounded below there; a lower bound on $\sitwo$ over an interval may
  certify ``no third eigenvalue hiding here''.
\item[A19. Eigenvector continuation.] Track $c(\lam)$ rather than $\sig(\lam)$. The coefficient
  vector swings as $\lam$ crosses an eigenvalue, and the swing is a large $O(1)$ signal where the
  tension change may be small. \textbf{Tested.}
\end{description}

\subsection{Family 6 --- model-based and statistical}
\begin{description}[leftmargin=1.2em,style=nextline,itemsep=0.25em]
\item[A20. Gaussian process with a level-set acquisition.] Evaluations cost 14--100\,ms, expensive
  enough that model-based placement can pay. Concern: a GP prior is stationary and $\sione$ is
  emphatically not --- deep narrow dips on a flat plateau. Would need the cone model as a mean
  function, which is A21.
\item[A21. Parametric cone-mixture fit.] Fit
  $\sione \approx \min_i \sqrt{\varepsilon_i^2 + c^2(\lam-\lam_i)^2}$ over a window, then place the
  next sample where adding one more cone would be least detectable.
\item[A22. Multistart local minimization with clustering.] No certificate at all, and the starts
  have to come from somewhere --- which is A1. \textbf{Rejected as a primary}, kept as a
  refinement stage.
\item[A23. Stochastic global search (annealing, evolutionary, CMA-ES).] \textbf{Rejected.} No
  completeness certificate, and clause~(1) is a multiset match. These answer a different question.
\end{description}

\subsection{Family 7 --- exploiting problem structure}
\begin{description}[leftmargin=1.2em,style=nextline,itemsep=0.25em]
\item[A24. A coarse FEM counting oracle.] Borrow $N(\lam)$ from a cheap low-order FEM inertia
  count as the completeness certificate, keeping MPS for accuracy. It supplies precisely the
  missing A13, and the Weyl audit is a weak surrogate for it --- Weyl is blind at the top of the
  range and needs five distinct levels; inertia is exact and local.
\item[A25. Symmetry decomposition into irreps.] For a \emph{declared} symmetry, each irrep has a
  simple spectrum and multiplicity becomes bookkeeping across irreps rather than a numerical test.
  Falsifiable on \file{plus\_shape} ($D_4$), where the two measured doubles should split into
  distinct irreps. Constraint: declared only, never detected (X8).
\item[A26. Two-level rank.] Low rank for the scan, high for refinement, one factorization. Already
  half-measured: rank 111 locates $\lam_1$ to $3.2\times10^{-12}$ against $1.4\times10^{-10}$ at
  rank 61.
\item[A27. Basis-adaptive escalation.] When the audit says the basis is binding
  ($\text{gap} < 20\,\varepsilon/c$), raise the basis rather than the sampling. Half-measured:
  \file{rect81}'s 4094 evaluations did not fix it (\S\ref{sec:localizer}).
\end{description}

\subsection{Triage}

\begin{table}[htbp]
\centering
\caption{Triage of the 27 candidates \ap{8}.}
\begin{tabular}{ll}
\toprule
tier & approaches \\
\midrule
\textsc{test}   & A2, A6, A7, A11, A14, A19 --- cheap, and each falsifies something load-bearing \\
\textsc{keep}   & A3, A4, A5, A8, A15, A24, A25, A26, A27 --- plausible, needs the tested ones first \\
\textsc{later}  & A12, A17, A18, A20, A21 --- real but expensive to evaluate \\
\textsc{reject} & A9, A13, A22, A23 --- reasons recorded above \\
\bottomrule
\end{tabular}
\end{table}

The six in \textsc{test} were chosen because each one, if it fails, removes a whole family: A2, if
stratification does not beat uniform, no sampling schedule will and the answer must come from the
certificate side; A6, if pilot secants underestimate the true slope, every Lipschitz certificate
in the design is unsound as built; A7, if Lipschitz-free subdivision matches it, the calibration
problem disappears; A11, if the phase of $\det A_B$ survives, we get an \emph{exact} count, which
nothing else offers; A14, if the derivative identity holds, four other approaches become available
for one gradient; A19, if the eigenvector swing is coarser-samplable than the dip, the stride
problem of \S\ref{sec:stride} is solved outright \ap{8}.

\section{The six tests}
\label{sec:tests}

All six were run (\file{notes/test\_approaches.py}). Table~\ref{tab:six} is the summary; each is
then given in full.

\begin{table}[htbp]
\centering
\caption{Outcomes of the six \textsc{test}-tier approaches \lf{22}.}
\label{tab:six}
\begin{tabular}{llp{7.2cm}}
\toprule
& outcome & evidence \\
\midrule
A2  & \textsc{falsified}      & worse or equal in 22 of 28 comparisons \\
A6  & \textsc{unsound}        & \file{GWW1} at $n_{\mathrm{pilot}}=20$ reads $0.672$: the
                                certificate is \emph{false}, not merely loose \\
A7  & \textsc{strong}         & mean recall $0.95$ against uniform's $0.75$, and no constant to
                                calibrate \\
A11 & \textsc{falsified}      & 455 spurious sign flips against 9 roots \\
A14 & \textsc{confirmed}      & $4.65\times10^{-9}$ once the $k$-dependent column scale is
                                corrected for \\
A19 & \textsc{reinterpreted}  & false as proposed; the peaks are crossings, not eigenvalues \\
\bottomrule
\end{tabular}
\end{table}

\subsection{A2 --- Weyl-stratified sampling: falsified}

\paragraph{What it is.} Sample equispaced in $N(\lam)$ rather than in $\lam$, so that points
concentrate where the spectral density is high.

\paragraph{The falsifiable claim.} At equal point count, stratified recall $\ge$ uniform recall,
with the gain largest where $\min\text{gap}/\text{mean gap}$ is smallest --- i.e.\ on
\file{mushroom} ($0.127$) and \file{H\_shape} ($0.210$).

\paragraph{The evidence.} Measured against uniform at equal point count, on all seven domains at
$n = 40, 60, 90, 130$: \textbf{28 comparisons, and stratification is worse or equal in 22 of
them.} At $n = 130$ it is worse on four of seven \lf{22}.

\paragraph{Why it kills the family.} The reason is structural. $N(\lam)$ is nearly \emph{linear}
in $\lam$, so stratifying by it is nearly uniform sampling, and the deviations it does introduce
are not aligned with where the local gaps are tight. \textbf{Local clustering is a fluctuation
phenomenon --- level repulsion, accidental near-degeneracy --- not a density phenomenon, and
Weyl's law describes only the density.} No a-priori schedule computed from Weyl can concentrate
samples where they are actually needed \lf{22}, \ap{9}. This is the formal reason F3 must correct
the stride \emph{in the loop}: the information needed to choose it in advance does not exist in
$N(\lam)$.

\subsection{A6 --- Lipschitz calibration: unsound at small pilots}

\paragraph{What it is.} Piyavskii--Shubert branch-and-bound with
$L = \gamma \max|d\sig/d\lam|$ estimated over $n_{\mathrm{pilot}}$ pilot points, $\gamma = 2$.
This is the originating proposal for the whole search.

\paragraph{The falsifiable claim.} The certificate is only as good as $L$, and $L$ from pilot
secants underestimates the true slope on some real domain.

\paragraph{The evidence.} Ratios of the pilot-derived $L$ to the true maximum slope from the fine
scan; values above 1 mean the certificate holds (Table~\ref{tab:a6}).

\begin{table}[htbp]
\centering
\caption{A6: pilot-derived $L$ against the true maximum slope \lf{22}.}
\label{tab:a6}
\begin{tabular}{lrrrrr}
\toprule
domain & true max & $n=20$ & $n=40$ & $n=80$ & $n=160$ \\
\midrule
\file{L\_shape} & $0.2013$ & $1.894$ & $1.989$ & $1.996$ & $2.000$ \\
\file{mushroom} & $0.3589$ & $1.524$ & $1.945$ & $1.974$ & $1.995$ \\
\file{GWW1}     & $0.7612$ & $\mathbf{0.672}$ & $1.974$ & $1.988$ & $1.999$ \\
\file{H\_shape} & $0.2438$ & $1.886$ & $1.941$ & $1.985$ & $1.991$ \\
\file{plus}     & $0.2738$ & $1.939$ & $1.986$ & $1.996$ & $1.999$ \\
\file{chevron}  & $0.0472$ & $1.993$ & $2.000$ & $2.000$ & $2.000$ \\
\file{rect81}   & $0.2264$ & $1.917$ & $1.945$ & $1.951$ & $1.961$ \\
\bottomrule
\end{tabular}
\end{table}

\textbf{\file{GWW1} at $n = 20$ reads $0.672$: the pilot underestimates the true slope even with
$\gamma = 2$, so the certificate is \emph{false}.} It is the domain with the densest spectrum, and
the pilot simply never lands on a steep wall \lf{22}.

\paragraph{Why this is \emph{unsound}, not merely loose.} The distinction matters and is worth
stating plainly. A loose constant costs evaluations; an \emph{underestimated} one voids the
guarantee. \file{search.md} \S6's conditional claim --- ``no $x$ with
$\sione(x) \le \tau_{\mathrm{scan}}$ outside the returned set, \emph{given} $L$'' --- is
conditional on something that can be false, and the failure mode is a silently missed eigenvalue
rather than wasted work \ap{9}. Everything converges to $2.000$ by $n = 40$, so the failure is not
asymptotic; it is that ``a fine enough pilot'' is domain-dependent and not knowable in advance,
which is \S\ref{sec:stride}'s stride problem wearing a different hat. A6 is usable only if
something else sizes the pilot --- which is why it survives as \textsc{wounded} rather than dead.

\subsection{A7 --- DIRECT, Lipschitz-free: strong}

\paragraph{What it is.} Keep every interval that is potentially optimal for \emph{some} Lipschitz
constant --- the lower-right convex hull of (width, value) --- plus every interval already below
$\tau$. There is no constant to calibrate, so A6's failure mode cannot occur.

\paragraph{The falsifiable claim.} Extended to keep every interval potentially containing a point
below $\tau$, DIRECT reaches recall $1.0$ at fewer evaluations than a uniform grid.

\paragraph{The evidence.} At about 230 evaluations (Table~\ref{tab:a7}).

\begin{table}[htbp]
\centering
\caption{A7: DIRECT against a uniform grid at the same evaluation count \lf{22}.}
\label{tab:a7}
\begin{tabular}{lrr}
\toprule
domain & DIRECT recall & uniform recall at the same $n$ \\
\midrule
\file{L\_shape} & $1.000$ & $0.778$ \\
\file{mushroom} & $1.000$ & $0.917$ \\
\file{GWW1}     & $0.800$ & $0.867$ \\
\file{H\_shape} & $0.875$ & $0.625$ \\
\file{plus}     & $1.000$ & $0.875$ \\
\file{chevron}  & $1.000$ & $0.333$ \\
\file{rect81}   & $1.000$ & $0.857$ \\
\bottomrule
\end{tabular}
\end{table}

\textbf{Six of seven at or above uniform, mean recall $0.95$ against $0.75$}, and it beats
\S\ref{sec:localizer}'s audit-guided loop on \file{mushroom}, \file{chevron} and \file{plus} while
costing $7\times$ less than it on \file{H\_shape} and $18\times$ less on \file{rect81} \lf{22}.

\paragraph{Caveat, recorded.} This simulation interpolates from the cached fine scan, so it does
not pay real evaluation cost and the interpolation smooths the curve. The structural result --- no
constant needed, and it beats uniform --- does not depend on that \lf{22}.

\subsection{\texorpdfstring{A11 --- the argument principle on $\det A_B$: falsified}{A11 -- the argument principle on det A\_B: falsified}}

\paragraph{What it is.} With a square collocation, $\det A_B(\lam)$ is analytic in $\lam$ and its
sign changes count the odd-multiplicity roots. This was the only exact-count candidate available
to MPS.

\paragraph{The falsifiable claim, and why it was not the usual objection.} The determinant method
is known to be numerically bad --- that is exactly why Betcke and Trefethen introduced the
subspace angle and why this project uses the GSVD. But that objection is about the
\emph{magnitude}, and counting needs only the \emph{sign}, which comes from LU pivot parity and is
a different question. The narrow claim was that the phase might survive where the magnitude does
not \ap{3}.

\paragraph{The evidence.} It is not a different question. On 900 samples:

\begin{table}[htbp]
\centering
\caption{A11: sign flips of $\det A_B$ against the true eigenvalue count \lf{22}.}
\begin{tabular}{lrrr}
\toprule
domain & $\log|\det|$ span & sign flips & true eigenvalues \\
\midrule
\file{L\_shape} & 214 orders & 464 (455 spurious) & 9 \\
\file{mushroom} & 163 orders & 456 (444 spurious) & 12 \\
\bottomrule
\end{tabular}
\end{table}

Every true eigenvalue does have a flip near it --- 9/9 and 12/12 --- but at 2\% precision that is
worthless.

\paragraph{The finding, which is stronger than the standard objection.}
\textbf{The conditioning destroys the determinant's \emph{topology}, not merely its magnitude}
\lf{22}. That is a strictly stronger statement than the usual complaint about determinant methods,
and it kills A11 and most of A12 (Beyn on a square MPS collocation) with it, because they rest on
the same matrix. With A13 already rejected, MPS is left with no exact counting oracle of its own,
which is what makes A24 (borrowing one from FEM) the standing option.

\subsection{\texorpdfstring{A14 --- the exact $k$-derivative: confirmed, after a correction}{A14 -- the exact k-derivative: confirmed, after a correction}}

\paragraph{What it is.} Every particular solution here is $\Psi(kx)$ about its own centre, so
Euler's identity for a homogeneous argument gives
\[
  \frac{\partial \psi}{\partial k} \;=\; \frac{x\cdot\nabla\psi}{k}
\]
exactly. $d\sione/d\lam$ then follows by first-order perturbation of the GSVD.

\paragraph{The falsifiable claim.} The analytic derivative matches a central finite difference to
the accuracy the finite difference deserves.

\paragraph{The evidence, including the initial failure.} Measured against a central difference:
\begin{center}
\begin{tabular}{ll}
\toprule
\verb|PlaneWave|   & $3.12\times10^{-10}$ \quad (finite-difference noise floor $2.23\times10^{-10}$) \\
\verb|FB @ corner| & $2.35\times10^{1}$ \;$\to$\; $4.65\times10^{-9}$ after the correction below \\
\bottomrule
\end{tabular}
\end{center}

\textbf{The raw identity fails for Fourier--Bessel}, and the reason is the $k$-dependent
\emph{column scale}. With $\psi_s = e^{-L(k)}\psi$ and $L$ spatially constant,
\[
  \frac{\partial \psi_s}{\partial k}
  \;=\; \frac{x\cdot\nabla\psi_s}{k} \;-\; L'(k)\,\psi_s ,
\]
and $|dL/dk|$ reaches $1.74$ for the corner family against exactly $0$ for plane waves --- which is
precisely why one passed and one did not. With the correction both sit at the finite-difference
floor \lf{22}. The episode is worth keeping in the record: the identity is exact mathematics, and
what broke it was an implementation detail of the equilibration, not the mathematics.

\paragraph{What it opens.} $d\sig/d\lam$ is available for one gradient evaluation plus a scalar
per column, rather than two extra tension evaluations. And $x\cdot\nabla\psi$ is the Rellich
multiplier, which the boundary-integral layer is already planned around. This opens A10 (a Hermite
envelope, far tighter than a secant bound), A15 (sign-change bracketing with no extra evaluations),
A16 (curvature for $\delta_{\mathrm{res}}$ directly), and a $c$ calibrated exactly rather than by
secants --- which is the direct repair for A6 \lf{22}, \ap{9}.

\subsection{A19 --- the eigenvector swing: falsified as proposed, strong as reinterpreted}
\label{sec:a19}

\paragraph{What it is, as proposed.} Track $c(\lam)$ rather than $\sig(\lam)$: the angle between
$c(\lam)$ and $c(\lam+h)$ was hypothesized to peak \emph{at} eigenvalues, giving an indicator that
survives a coarser stride than $\sione$'s dip.

\paragraph{The evidence against the original hypothesis.} \textbf{It does not: recall $0.333$ at
4 samples per gap, falling to $0.000$ at 8 and finer} \lf{22}. The finer the stride, the worse it
gets --- the signature of an indicator that is firing somewhere other than where it was expected.

\paragraph{The reinterpretation.} The swing peaks where the \emph{minimising vector changes
identity}, which is at a branch \emph{crossing} --- and crossings lie \emph{between} eigenvalues.
This is far more useful than the original hypothesis: a peak in every gap partitions $[a,b]$ into
cells holding exactly one eigenvalue each, which is a completeness certificate of a kind no
sampling schedule gives. Measured as ``gaps hit'' out of the gaps between known dips
(Table~\ref{tab:a19}).

\begin{table}[htbp]
\centering
\caption{A19b: the swing as a cell decomposition \lf{22}. Timing is per sample, for
\texttt{coef}.}
\label{tab:a19}
\begin{tabular}{lllllrr}
\toprule
domain & 3/gap & 4/gap & 6/gap & 8/gap & extra peaks at 4/gap & ms per sample \\
\midrule
\file{L\_shape} & 8/8   & 8/8   & 8/8   & 9 (1 extra) & 0 & $13.7$ \\
\file{mushroom} & 10/11 & 11/11 & 11/11 & 11/11       & 0 & $32.7$ \\
\file{GWW1}     & 14/14 & 14/14 & 14/14 & 14/14       & 0 & $111.3$ \\
\file{H\_shape} & 7/7   & 7/7   & 7/7   & 7/7         & 0 & $81.7$ \\
\file{plus}     & 7/7   & 7/7   & 7/7   & 7/7         & 0 & $78.9$ \\
\file{chevron}  & 2/2   & 2/2   & 2/2   & 2/2         & 0 & $106.9$ \\
\file{rect81}   & 12/13 & 12/13 & 13/13 & 13/13       & 0 & $10.7$ \\
\bottomrule
\end{tabular}
\end{table}

\textbf{At four samples per Weyl gap, six of seven domains give a perfect cell decomposition},
and \file{rect81} needs six. Against the 18 per median gap that $\sione$ dip detection needs for
full recall (\S\ref{sec:stride}), that is a $3\times$ to $4.5\times$ reduction in samples.

\textbf{And it is nearly free.} \verb|coef| costs $13.7$\,ms against \verb|tensions|' $13.6$ on
\file{L\_shape}, $111.3$ against $104.8$ on \file{GWW1}, $81.7$ against $76.6$ on
\file{H\_shape} --- \textbf{5 to 7 percent} --- because both need the same regularization and
factorization, and the \verb|separate|-mode GSVD is comparable to \verb|gsvdvals| \lf{22}.

Extra peaks are 0 or 1--2 and are the safe direction: an extra cell is an empty cell, which is
cheap to detect and discard, whereas a missing cell is a missing eigenvalue.

\paragraph{Three caveats, all real and all recorded.} The denominator is gaps between
\emph{detected} dips, so this measures ``does it separate the levels we already know about'', not
``does it find a hidden pair'' --- the pair-resolution question is untouched. The
$5\times\text{median}$ peak threshold is a parameter and has not had the margin study that $G$
and $\tau_{\mathrm{gram}}$ got. And it \emph{brackets} rather than locates: refinement still runs
inside each cell, though refinement inside a bracket known to hold exactly one minimum is the easy
and robust case, which is the entire point \lf{22}.

\subsection{Shortlist after the six tests}

\begin{center}
\begin{tabular}{ll}
\toprule
\textsc{keep, tested}   & A7 (DIRECT), A19b (swing cell decomposition), A14 (exact derivative) \\
\textsc{keep, untested} & A3, A5, A8, A15, A24, A25, A26, A27 \\
\textsc{dead}           & A2, A11, A12-square, A9, A13, A22, A23 \\
\textsc{wounded}        & A6 --- unsound at small pilots; usable only with A7 or A19b sizing it \\
\bottomrule
\end{tabular}
\end{center}

The two survivors attack the same defect from opposite sides. \textbf{A19b gives the partition;
A7 gives the search inside it, with no constant to calibrate.} They compose: use the swing at 4
per gap to cut $[a,b]$ into one-eigenvalue cells, then run DIRECT inside each cell against $\tau$,
then refine. That combination has a completeness argument that neither the Lipschitz certificate
(A6, measured unsound) nor the Weyl audit (\S\ref{sec:weyl}, blind at the top and needing five
levels) can provide, and it is the first candidate in this project that does \lf{22}.

\section{The swing on a pair: a detection result}
\label{sec:swingpairs}

\S\ref{sec:a19} left A19b's central question open: its cells were measured against gaps between
\emph{detected} dips, so it had never been shown to split a pair that $\sione$ presents as one
dip. That is the failure mode that \lf{9.4} and \lf{19} call real, and it decides whether the
composed method is a capability win or only a cost win. The instrument is
\verb|rect(1+delta, 1)| with modes $(2,1)$ and $(1,2)$, in one window holding the pair and no
neighbour (\file{notes/swing\_pairs.py}).

\begin{table}[htbp]
\centering
\caption{Pair detection at stride $h = 1.005$, about 12 samples per gap \lf{23}.}
\label{tab:swingpair}
\begin{tabular}{lccl}
\toprule
$d(\lam)$ & $\sione$ dips & swing peaks & verdict \\
\midrule
$5.139\times10^{0}$  & 2 & 1 & $\sione$ already splits \\
$1.700\times10^{0}$  & 2 & 1 & $\sione$ already splits \\
$5.834\times10^{-1}$ & 1 & 1 & swing sees it, dip does not \\
$1.769\times10^{-1}$ & 1 & 1 & swing sees it \\
$5.913\times10^{-2}$ & 1 & 1 & swing sees it \\
$1.776\times10^{-2}$ & 1 & 1 & swing sees it \\
$5.921\times10^{-3}$ & 1 & 1 & swing sees it \\
$1.776\times10^{-3}$ & 1 & 0 & neither \\
$0$ (exact)          & 1 & 0 & degenerate control \\
\bottomrule
\end{tabular}
\end{table}

At that stride $\sione$ stops splitting the pair below $d \sim 1.7$, and the swing keeps splitting
it down to $d \sim 5.9\times10^{-3}$: \textbf{a factor of about 290}, at the same sampling cost.
In stride units the swing fires at $d/h = 5.9\times10^{-3}$ --- it resolves a pair about
$170\times$ finer than the sample spacing \lf{23}.

\paragraph{This contradicts a recorded conclusion.} \file{search.md} \S3 says ``there is no
detector here that sees a pair before $\sione$ does'', and \lf{11} measured the offset detector
firing ``at exactly the same coarseness at which $\sione$ first shows two minima''. Both are
correct \emph{about the offset}, and both were generalized into a claim about detectors in
general. The swing is a counterexample, and it works because it keys on something neither of those
does: \textbf{the identity of the minimising eigenfunction, not the shape of the tension curve.}
The two eigenfunctions stay near-orthogonal at $|\cos| \sim 1/\sqrt{n_{\mathrm{int}}}$ far past
where the tension merges them, and that orthogonality is what survives \lf{23}.

\paragraph{The degenerate control is the other half of the result.} At $\delta = 0$ the swing
reads \emph{zero} peaks. That is correct and not a failure: a true double has no crossing, because
there is no other branch for the minimiser to switch to. So
\begin{center}
\begin{tabular}{l@{\ }l}
\toprule
swing peak inside a dip & $\to$ two \emph{distinct} eigenvalues \\
no swing peak           & $\to$ a genuine degeneracy, \emph{or} a pair below the swing's own limit \\
\bottomrule
\end{tabular}
\end{center}
The first arrow is a multiplicity discriminator of a kind (C1)$+$(C2) does not have: the
double-count case (the gap rule reporting multiplicity 2 at a \emph{simple} eigenvalue once
$d \lesssim \varepsilon/c$) is exactly where the swing says ``distinct'', whereas (C2) currently
has to reject it by a monotonicity argument on $\sitwo$. The second arrow is still ambiguous, so
this does not remove the resolution limit; \textbf{it moves it} \lf{23}.

\paragraph{Where the limit sits, and it is not the pencil.}
\begin{center}
\begin{tabular}{ll}
\toprule
$h = 3.770$ & fires to $d = 5.9\times10^{-2}$, not $1.8\times10^{-2}$ \\
$h = 1.885$ & fires to $d = 5.9\times10^{-2}$, not $1.8\times10^{-2}$ \\
$h = 1.005$ & fires to $d = 5.9\times10^{-3}$, not $1.8\times10^{-3}$ \\
\bottomrule
\end{tabular}
\end{center}
$d_{\min}$ tracks $h$, roughly $h/30$ to $h/170$. It is \emph{not} $\varepsilon/c$, which is about
$3.5\times10^{-9}$ here --- six orders below where the swing stops. \textbf{So the swing's limit is
set by the stride and the peak threshold, not by the basis}, and there is a great deal of headroom
left: a finer stride buys proportionally finer pair detection, which is not true of any dip-based
detector \lf{23}.

\paragraph{Caveats.} One domain, one mode pair, one basis, $\text{thresh} = 5\times\text{median}$.
The threshold has had no margin study --- $G$ and $\tau_{\mathrm{gram}}$ both turned out
non-delicate and this one has not been checked at all. And $d_{\min} \sim h/30$ to $h/170$ is a
$5\times$ spread across three strides, so the scaling is indicative rather than established
\lf{23}.

\section{The composed candidate, scored}
\label{sec:composed}

\verb|tests/searches.py::swing_direct_search|: swing cells at 6 samples per estimated gap, DIRECT
inside each cell against $\tau = 0.5\,S$, Brent in the resulting one-root bracket, (C1)$+$(C2) for
multiplicity. Scored through the harness on five real domains with \verb|reference| table truth,
windows cut inside a tabulated gap so the truth is \emph{complete}:

\begin{center}
\begin{tabular}{ll@{\ }r@{\qquad}ll@{\ }r}
\toprule
\file{L\_shape} & $[4.82, 49.35]$ & 9 & \file{GWW1}    & $[1.27, 11.94]$   & 8 \\
\file{mushroom} & $[2.75, 33.39]$ & 8 & \file{H\_shape}& $[3.87, 25.58]$   & 8 \\
\file{chevron}  & $[19.83, 192.25]$ & 8 \\
\bottomrule
\end{tabular}
\end{center}

\begin{table}[htbp]
\centering
\caption{End-to-end scores on the five-domain roster \lf{24}.}
\label{tab:e2e}
\begin{tabular}{lrrrr}
\toprule
method & mean recall & total evals & spurious & mult errors \\
\midrule
\verb|uniform_grid|   & $0.881$ & 1010 & 0 & 0 \\
\verb|lappy_repaired| & $0.319$ & 1946 & 0 & 0 \\
swing $+$ DIRECT      & $0.975$ & 2221 & 0 & 0 \\
\bottomrule
\end{tabular}
\end{table}

\begin{table}[htbp]
\centering
\caption{Per case, swing $+$ DIRECT \lf{24}. The exclamation mark on \file{H\_shape} is the Weyl
deficit exceeding its $0.5$ threshold --- the flag firing on the one miss.}
\label{tab:percase}
\begin{tabular}{lrrrrrrrr}
\toprule
case & recall & prec & miss & spur & multX & max err & evals & weyl \\
\midrule
\file{L\_shape} & $1.000$ & $1.000$ & 0 & 0 & 0 & $2.3\times10^{-8}$ & 473 & $0.41$ \\
\file{mushroom} & $1.000$ & $1.000$ & 0 & 0 & 0 & $4.1\times10^{-9}$ & 442 & $0.02$ \\
\file{GWW1}     & $1.000$ & $1.000$ & 0 & 0 & 0 & $8.8\times10^{-9}$ & 454 & $0.33$ \\
\file{H\_shape} & $0.875$ & $1.000$ & 1 & 0 & 0 & $8.0\times10^{-9}$ & 385 & $0.76$\;! \\
\file{chevron}  & $1.000$ & $1.000$ & 0 & 0 & 0 & $9.2\times10^{-9}$ & 467 & $0.12$ \\
\bottomrule
\end{tabular}
\end{table}

\textbf{Four of five perfect, no spurious roots and no multiplicity errors anywhere}, and the one
miss is \emph{flagged} --- \file{H\_shape} reads $0.76$ against the $0.5$ threshold. That is
clause~(3) of the operational definition doing its job: the failure is reported rather than
silent, which is the difference between recoverable and fatal inside a shape-optimization loop
\lf{24}.

\subsection{The honest ledger}

\paragraph{It costs $2.2\times$ the uniform grid} for $+0.094$ recall. Recall is the headline
metric by design --- a missed eigenvalue cannot be repaired downstream --- so this is the right
trade under the stated principle, but \textbf{it is a trade and not a free win} \lf{24}.

\paragraph{\texttt{lappy\_repaired}'s $0.319$ is not a fair comparison.} It is starved by the
harness's crude \verb|_spacing| estimator, exactly the caveat \lf{11} records: a coarse pass sees
a pair as one dip, concludes the interval holds one eigenvalue, and sizes a 20-point grid.
\textbf{The number measures the estimator, not lappy's architecture.} Its own recall table, with
\verb|n_pts| set explicitly, is 10/12 \lf{24}.

\paragraph{The accuracy column is F11 working, not a defect.} swing$+$DIRECT reports maximum errors
of $4\times10^{-9}$ to $2\times10^{-8}$ against a requested $\verb|prec| = 10^{-8}$; the uniform
grid reports $8\times10^{-10}$ to $6\times10^{-9}$ because it refines to
$\verb|xatol| = 10^{-13}$ regardless of what was asked. The composed search stops at
$\max(\verb|prec|\cdot\lam,\ \delta_{\mathrm{res}})$ and the first term binds. Doing better than
asked is allowed; doing it at $2\times$ the refinement cost on every root is what \sm{2.6} rules
out \lf{24}.

\subsection{Two efficiency findings}

The first version re-scanned the \emph{cell} grid to estimate the plateau, having already scanned
a 64-point pilot for the spacing. The plateau is a median over an $O(1)$-varying curve and needs no
cell-stride resolution; folding both into one 48-point pilot cut 2535 evaluations to 2221 with
identical recall.

Dropping to 4 samples per gap and a DIRECT budget of 10 cuts cost further to 1583 --- $1.6\times$
--- but recall falls to $0.925$ (\file{H\_shape} $0.875 \to 0.750$, \file{chevron}
$1.000 \to 0.875$). \S\ref{sec:a19} measured 4 per gap as sufficient for the \emph{cells}, and it
is; \textbf{what fails at 4 is the DIRECT search inside them}. The two budgets are separable and
should be tuned separately \lf{24}.

\subsection{What the roster scoring left untested}

Five domains, all with table truth, all at low $\lam$. No pair below the swing's own limit appears
in any of these windows, so \S\ref{sec:swingpairs}'s $290\times$ is not exercised here --- the
composed search had not yet been scored on a case where it should beat a dip-based method
outright. \verb|rect(1+delta, 1)| is the instrument and it was not wired into the harness
\lf{24}. That is what \S\ref{sec:sep} does, and what it found was a correction.

\section{Detection is not separation: a correction to the pair result}
\label{sec:sep}

\verb|rect(1+delta, 1)| was wired in as harness cases: window $[10, 88]$, holding $(1,1)$, the
$(2,1)/(1,2)$ pair and $(2,2)$; truth from the closed form and \emph{complete}; $\delta$ swept so
the pair separation runs $5.8\times10^{-1}$ down to $0$.

\textbf{The first run failed, at every separation including the widest.} All three methods scored
recall $0.750$ --- three of four --- and the degenerate $\delta = 0$ case scored $1.000$. A window
whose \emph{split} cases fail while its \emph{degenerate} case passes says the pair is being
merged, not that a level is missing \lf{25}.

\subsection{The defect}

Instrumenting the $\delta = 10^{-2}$ case, $d = 5.83\times10^{-1}$:
\begin{center}
\begin{tabular}{ll}
\toprule
pilot sees 3 wells in a 4-eigenvalue window & $\to$ estimated gap $26.0$ \\
cell grid $n = 19$ over $[10, 88]$          & $\to$ stride $h = 4.33$ \\
crossing truly at $48.86$                   & \\
swing peak \emph{reported} at $46.83$       & \\
cells $[10,\ 29.5,\ 46.83,\ 64.17,\ 88]$    & $\to$ \textbf{both pair members in one cell} \\
\bottomrule
\end{tabular}
\end{center}
The swing \emph{detected} the crossing. It localized it only to $\pm h/2$, and with $d < h$ the
cell boundary landed outside the pair instead of between its members.

\paragraph{The self-correction, stated plainly.}
\textbf{Section~\ref{sec:swingpairs} measured detection, not separation, and one was generalized
into the other.} Its in-pair test counted a peak as in-pair if it lay within
$[\lam_a - h,\ \lam_b + h]$ --- a tolerance of a full stride on each side. Every number in
\S\ref{sec:swingpairs} stands as a statement about \emph{detection}; none of them was ever a
statement about splitting. \textbf{A pair detector and a pair separator are different things, and
the $290\times$ is the first, not the second} \lf{25}.

\subsection{The repair: a peak is a feature to bisect toward}

The crossing lies in whichever half of an interval carries the larger angle; the halves are equal
width so the comparison is fair; and each level costs one mode evaluation. Cost is
$\log_2(h/d)$ --- the same logarithmic behaviour adaptive subdivision has on the tension, obtained
here from the eigenvector instead --- and it converges to the crossing \emph{regardless of how
coarse the initial stride was}. That last property matters, because the stride here was set by a
pilot that merged the pair and therefore could not have been chosen correctly in advance \lf{25}.

\begin{table}[htbp]
\centering
\caption{swing$+$DIRECT on \texttt{rect(1+delta, 1)}, before and after crossing localization
\lf{25}.}
\label{tab:rectpair}
\begin{tabular}{lrrrr}
\toprule
$d(\lam)$ & recall before & recall after & \verb|uniform_grid| & \verb|lappy_repaired| \\
\midrule
$5.834\times10^{-1}$ & $0.750$ & $1.000$ & $0.750$ & $0.250$ \\
$1.769\times10^{-1}$ & $0.750$ & $1.000$ & $0.750$ & $0.750$ \\
$5.913\times10^{-2}$ & $0.750$ & $1.000$ & $0.750$ & $0.250$ \\
$1.776\times10^{-2}$ & $0.750$ & $1.000$ & $0.750$ & $0.250$ \\
$5.921\times10^{-3}$ & $0.750$ & $0.750$ & $0.750$ & $0.750$ \\
$1.776\times10^{-3}$ & $0.750$ & $0.750$ & $0.750$ & $0.500$ \\
$5.922\times10^{-4}$ & $0.750$ & $0.750$ & $0.750$ & $0.500$ \\
$0$ (exact double)   & $1.000$ & $1.000$ & $1.000$ & $0.250$ \\
\midrule
mean                 & $0.781$ & $0.906$ & $0.781$ & $0.438$ \\
\bottomrule
\end{tabular}
\end{table}

\textbf{The uniform grid never splits this pair at any separation tested, including the widest},
because its stride comes from the same merged-well estimate. The composed search splits it down to
$d = 1.78\times10^{-2}$, so its advantage over the control is at least $32\times$ in separation,
and the true factor is larger since the control's own limit is above $5.8\times10^{-1}$.

The roster of \S\ref{sec:composed} is unchanged at mean recall $0.975$, at 2596 evaluations
against 2221 --- localization costs about 17\% and buys the pair regime \lf{25}.

\subsection{The remaining frontier}

At $d \le 5.9\times10^{-3}$ the composed search still scores $0.750$ with one spurious root and one
multiplicity error. The cell boundary \emph{is} found (\S\ref{sec:swingpairs} says the swing
detects to $5.9\times10^{-3}$), but the two refinements in adjacent cells converge to the same
minimum. The tension floor at this pair is $9.1\times10^{-8}$ to $1.1\times10^{-7}$, so
$\varepsilon/c$ is about $2\times10^{-6}$ and $d = 5.9\times10^{-3}$ is $3000\times$ above it:
\textbf{the pair is resolvable in principle and this is a search failure, not a basis limit.} Two
candidates, neither tested \lf{25}:
\begin{enumerate}[nosep]
\item DIRECT's fixed budget is too small for a cell whose dip sits near its edge;
\item the Gram-angle collapse phase is \emph{not implemented} in \verb|swing_direct_search|, so a
  duplicate arising from two adjacent cells is reported rather than collapsed and recounted.
\end{enumerate}
The second is a direct instance of the interface finding of \S\ref{sec:interface}: the collapse
decision is exactly the one the harness could not see before \verb|modes(x)| was available.

\section{Open questions and what is untested}
\label{sec:open}

Collected from the ``what was not checked'', ``still untested'' and ``remaining frontier''
passages of the sources.

\paragraph{On the stride rule.} The 18-per-gap number is a \emph{fitted} value on seven domains,
not a bound. The governing quantity --- the minimum local gap --- is not predictable a priori, so
this can only ever be a starting point for F3's in-loop correction \sfe{6}.

\paragraph{On the accuracy contract.} $C$ in the accuracy clause remains uncalibrated; clause~(3)
is the workaround. And the wide bracket used for (C2) was taken as $0.02\lam$, an ad-hoc fraction;
it should be derived from the refinement's own parabola, and is not \sfe{6}.

\paragraph{On the roster's coverage.} Everything above is at low $\lam$ --- the roster tops out at
90 and mostly sits at 30--40. There is no domain with a hole, and no curved boundary except
\file{cut\_square}, which was dropped for a basis-planning reason rather than a curve one
\sfe{6}, \lf{20}. The composed candidate's five-domain scoring is likewise all low $\lam$ and all
table truth \lf{24}.

\paragraph{On the swing.} One domain, one mode pair, one basis for the pair study. The
$5\times\text{median}$ peak threshold has had no margin study, unlike $G$ and
$\tau_{\mathrm{gram}}$. The $d_{\min} \sim h/30$ to $h/170$ scaling is indicative, not established
--- a $5\times$ spread over three strides \lf{23}.

\paragraph{On DIRECT.} The A7 numbers come from a simulation that interpolates the cached fine
scan, so it pays no real evaluation cost and the interpolation smooths the curve \lf{22}.

\paragraph{On the sub-$5.9\times10^{-3}$ pair regime.} Both candidate explanations --- an
under-sized DIRECT budget in an edge-dip cell, and the missing Gram-angle collapse --- are
untested \lf{25}.

\paragraph{Structurally unbuilt.} F6 (two-level rank from one factorization) is measured and not
built. The at-level prong of the deficit response and the $20\,\varepsilon/c$ termination guard
(\S\ref{sec:localizer}) are both specified and neither is implemented. The interface widening of
\S\ref{sec:interface} gates everything about multiplicity.

\paragraph{Still open in the enumeration.} A24 (a coarse FEM inertia count as a counting oracle) is
the standing answer to the fact that A11 and A13 are both dead and MPS therefore has no exact
count of its own. A25 (declared-symmetry irrep decomposition) would remove the (C1)/(C2) test
entirely for symmetric domains --- the class the solver will meet at the optimizers of a shape
optimization --- and is untested. A8 (the cone-model envelope), A15 (sign bracketing, now
available since A14 is confirmed), A26 and A27 are all in the \textsc{keep, untested} tier.

\section*{Note on two inconsistencies in the sources}
\addcontentsline{toc}{section}{Note on two inconsistencies in the sources}

Two points were left as the sources have them rather than silently resolved.

First, \file{docs/design/search.md} \S2.6 still states the scan stride as ``about six samples per
eigenvalue gap'' as part of its ``not more work than asked'' argument, while
\file{landscape\_findings} \S20.4 explicitly says that rule ``is wrong and should be replaced''.
The design document has not been updated. The clause it is attached to --- completeness is not a
precision question, and the stride does not tighten with \verb|prec| --- is unaffected by the
correction.

Second, the anti-feature list in \file{solver\_features.md} \S3 uses the labels A1--A9 for a set
of things measured not to work, while \file{approaches.md} uses A1--A27 for the candidate methods.
The two numberings are unrelated and collide (for instance, anti-feature A2 forbids a
$\sitwo$-value gate, while approach A2 is Weyl-stratified sampling). The anti-features were
relabelled X1--X9 in \S\ref{sec:anti} of this report; the underlying notes should probably adopt a
distinct prefix.

\end{document}
